\documentclass[a4paper,12pt]{article}

\usepackage[french]{babel}
\usepackage[T1]{fontenc} 
\usepackage[latin1]{inputenc}
%\usepackage{amsmath}

\author{R�mi Synave, Stefka Gueorguieva, Pascal Desbarats}
\title{Manuel de la biblioth�que POINT}
\date{}

\begin{document}
\maketitle

\section{Utilisation}
Cette biblioth�que impl�mente les structures de donn�es point 2D et point 3D permettant la manipulation de points du plan ou de l'espace.\\
Attention la structure \textit{point3d} n'est pas la m�me que la structure utilis�e pour le stockage des sommets d'un maillage. Pour les sommets d'un maillage, se reporter � la biblioth�que \textit{vertex.so}.\\

\section{Structures de donn�es}

Le point 2D est une structure comportant deux r��ls : sa composante en X et en Y. Le point 3D en comporte trois : composante en X,Y et Z.\\

\begin{verbatim}
typedef struct
{
  double x,y;
}point2d;
\end{verbatim}

\begin{verbatim}
typedef struct
{
  double x,y,z;
}point3d;
\end{verbatim}

\section{Fonctions}
\subsection{Fonctions utilisant \texttt{point2d}}

\textbullet \texttt{void point2d\_init(point2d *p, double x, double y)}\\
Initialisation d'un \texttt{point2d} avec deux r��ls.\\
~\\
\underline{Param�tres et type de retour :}\\
\texttt{p} : le point.\\
\texttt{x} : coordonn�es x.\\
\texttt{y} : coordonn�es y.\\
\texttt{retour} : aucun.\\


\textbullet \texttt{void point2d\_display(point2d p)}\\
Affichage d'un \texttt{point2d}.\\
~\\
\underline{Param�tres et type de retour :}\\
\texttt{p} : le point.\\
\texttt{retour} : aucun.\\


\subsection{Fonctions utilisant \texttt{point3d}}

\textbullet \texttt{void point3d\_init(point3d *p, double x, double y, double z)}\\
Initialisation d'un \texttt{point3d} avec trois r��ls.\\
~\\
\underline{Param�tres et type de retour :}\\
\texttt{p} : le point.\\
\texttt{x} : coordonn�es x.\\
\texttt{y} : coordonn�es y.\\
\texttt{z} : coordonn�es z.\\
\texttt{retour} : aucun.\\


\textbullet \texttt{void point3d\_display(point3d p)}\\
Affichage d'un \texttt{point3d}.\\
~\\
\underline{Param�tres et type de retour :}\\
\texttt{p} : le point.\\
\texttt{retour} : aucun.\\



\textbullet \texttt{int point3d\_equal(point3d p1, point3d p2)}\\
Test de l'�galite de deux \texttt{point3d}. Les deux points sont �gaux si leurs composantes sont �gales.\\
~\\
\underline{Param�tres et type de retour :}\\
\texttt{p1} : le premier point.\\
\texttt{p2} : le second point.\\
\texttt{retour} : aucun.\\

\end{document}
